\documentclass[12pt]{article}
\usepackage{amsmath,amssymb,amsfonts}
\usepackage[margin=1in]{geometry}

\begin{document}

\begin{center}
{\Large\bfseries Chapter 10, Problem 4}
\end{center}

\noindent\textbf{Problem.} Determine all the elements of the group of rotational symmetries of the regular tetrahedron (i.e., a tetrahedron whose sides are equilateral triangles). Put these elements into the classes to which they belong. Find all the subgroups. Are there any invariant subgroups?

\bigskip
\noindent\textbf{Solution.}

Label the four vertices of the regular tetrahedron as 1, 2, 3, 4. The rotation group of the regular tetrahedron is isomorphic to the alternating group $A_4$ (the group of even permutations of four objects), and has order 12.

\medskip
\noindent\textbf{The 12 elements:}

\begin{enumerate}
\item \textbf{Identity} (1 element): $e$

\item \textbf{Rotations by $120^\circ$ and $240^\circ$ about axes through each vertex and the center of the opposite face} (8 elements): There are 4 such axes, each giving two nontrivial rotations.
\begin{align*}
&\text{Axis through vertex 1: } (234), \; (243) \\
&\text{Axis through vertex 2: } (134), \; (143) \\
&\text{Axis through vertex 3: } (124), \; (142) \\
&\text{Axis through vertex 4: } (123), \; (132)
\end{align*}

\item \textbf{Rotations by $180^\circ$ about axes through midpoints of opposite edges} (3 elements):
\begin{align*}
&\text{Axis through edges 12 and 34: } (12)(34) \\
&\text{Axis through edges 13 and 24: } (13)(24) \\
&\text{Axis through edges 14 and 23: } (14)(23)
\end{align*}
\end{enumerate}

Total: $1 + 8 + 3 = 12$ elements. \checkmark

\medskip
\noindent\textbf{Conjugacy Classes:}

Elements are conjugate if and only if they have the same cycle structure in $A_4$:
\begin{align*}
&\mathcal{C}_1 = \{e\} \quad (1 \text{ element}) \\
&\mathcal{C}_2 = \{(123), (132), (124), (142), (134), (143), (234), (243)\} \quad (8 \text{ elements})
\end{align*}

Wait---in $A_4$ the 3-cycles actually split into two classes. Let us verify. In $S_4$, all 3-cycles are conjugate, but in $A_4$ they can split. We check: $(12)(34) \cdot (123) \cdot (12)(34) = (243)$. Since $(12)(34) \in A_4$, we get $(123) \sim (243)$. Also $(123)^2 = (132)$, and conjugating $(123)$ by $(13)(24)$ gives $(142)$.

In fact in $A_4$, the 3-cycles split into two conjugacy classes of 4 elements each:
\begin{align*}
&\mathcal{C}_2 = \{(123), (142), (134), (243)\} \quad (4 \text{ elements}) \\
&\mathcal{C}_3 = \{(132), (143), (124), (234)\} \quad (4 \text{ elements}) \\
&\mathcal{C}_4 = \{(12)(34), (13)(24), (14)(23)\} \quad (3 \text{ elements})
\end{align*}

Check: $1 + 4 + 4 + 3 = 12$. \checkmark

So there are \textbf{4 conjugacy classes}.

\medskip
\noindent\textbf{Subgroups:}

By Lagrange's theorem, subgroup orders must divide 12, so possible orders are 1, 2, 3, 4, 6, 12.

\begin{itemize}
\item \textbf{Order 1:} $\{e\}$

\item \textbf{Order 2:} $\{e, (12)(34)\}$, $\{e, (13)(24)\}$, $\{e, (14)(23)\}$ --- three subgroups.

\item \textbf{Order 3:} $\{e, (123), (132)\}$, $\{e, (124), (142)\}$, $\{e, (134), (143)\}$, $\{e, (234), (243)\}$ --- four subgroups (cyclic, generated by 3-cycles).

\item \textbf{Order 4:} $V = \{e, (12)(34), (13)(24), (14)(23)\}$ --- the Klein four-group. This is the only subgroup of order 4.

\item \textbf{Order 6:} There is no subgroup of order 6 in $A_4$. (If there were, it would be normal of index 2, but one can verify no such subgroup exists.)

\item \textbf{Order 12:} $A_4$ itself.
\end{itemize}

\medskip
\noindent\textbf{Invariant (Normal) Subgroups:}

A normal subgroup must be a union of conjugacy classes. The Klein four-group $V = \{e, (12)(34), (13)(24), (14)(23)\}$ is the union of $\mathcal{C}_1$ and $\mathcal{C}_4$, with $|V| = 1 + 3 = 4$, which divides 12. One can verify it is closed under multiplication. Therefore $V$ is a normal subgroup.

No other nontrivial proper normal subgroup exists in $A_4$. So the \textbf{only invariant subgroup} is the Klein four-group $V \cong \mathbb{Z}_2 \times \mathbb{Z}_2$.

The factor group $A_4/V$ has order $12/4 = 3$ and is therefore isomorphic to $\mathbb{Z}_3$.

\end{document}
